\chapter{Введение}
\label{chap:intro}

При гидродинамическом моделировании нефтяных пластов приходится искать компромисс между точностью и вычислительными затратами: детальная сетка отслеживает фронты вытеснения, но требует больших ресурсов, а укрупнение сетки ускоряет расчет ценой размытия фронта и завышения прогноза нефтедобычи. Классические методы интерполяции здесь бессильны: работая как фильтры нижних частот, они сглаживают картинку, не восстанавливая утраченные высокочастотные детали.

Цель работы — создать программный комплекс на базе нейросетевых подходов для быстрого и точного восстановления нестационарных полей давления и насыщенности по расчетам на грубой сетке. Подход использует суперразрешение на сверточных сетях, обученных на парах согласованных данных: быстрых расчетах на грубой сетке и эталонных — на детальной, полученных моделью трещиновато-пористых коллекторов.